\section{Trace-transfer reduction and the search domain}
\label{sec:transfer}

\subsection{Trace one and rank-one idempotence}

Use row vectorization,
\begin{align}
	\operatorname{vec}(Y) & =\sum_{u,u'}Y_{uu'}\ket{u,u'}.
\end{align}
The trace matrix of the letters is
\begin{align}
	T_{\alpha\beta} & =\tr M^{\alpha\beta},
	\nonumber                                   \\
	T               & =A_0\otimes\overline{A_0}
	+A_1\otimes\overline{A_1}
	+C_0\otimes\overline{C_0}
	+C_1\otimes\overline{C_1}.
	\label{eq:trace-matrix}
\end{align}
It represents the CP map
\begin{align}
	\mc E(Y)
	 & =A_0YA_0^\dagger+A_1YA_1^\dagger
	+C_0YC_0^\dagger+C_1YC_1^\dagger.
	\label{eq:transfer-map}
\end{align}
Indeed,
\begin{align}
	(KYK^\dagger)_{uu'}
	 & =\sum_{v,v'}K_{uv}Y_{vv'}\overline{K_{u'v'}},
	\nonumber                                        \\
	\operatorname{vec}(KYK^\dagger)
	 & =(K\otimes\overline K)\operatorname{vec}(Y).
\end{align}

Trace preservation of a refinement map in \cref{eq:rfp} requires
\begin{align}
	T_{\alpha\beta}
	 & =\tr G^{\alpha\beta}
	=\sum_\gamma
	\tr M^{\alpha\gamma}\tr M^{\gamma\beta}
	=(T^2)_{\alpha\beta}.
	\label{eq:idempotence}
\end{align}
Thus the bounded search imposes $T^2=T$. For an idempotent matrix,
\begin{align}
	\tr\rho_N(M) & =\tr(T^N)=\tr T=\rank T.
	\label{eq:trace-one}
\end{align}
The additional trace-one hypothesis therefore imposes $\rank T=1$.
This is a restriction on the physical tensor, not a normalization
convention. Replacing a nonzero tensor $M$ by $sM$, with $s>0$,
would require
\begin{align}
	(sT)^2=sT
	\quad\Longrightarrow\quad
	s^2T=sT
	\quad\Longrightarrow\quad
	s=1.
	\label{eq:no-rescaling}
\end{align}

Every nonzero positive map with one-dimensional range has the form
\begin{align}
	\mc E(Y) & =R\tr(LY),
	         &
	R,L      & \geq0.
	\label{eq:rank-one-map}
\end{align}
To see this, choose a positive input with nonzero positive output
$R$. Positive inputs span the Hermitian matrices, and the complex
range is one-dimensional, so every output is a scalar multiple of
$R$. The scalar functional is positive. Trace duality represents it
as $Y\mapsto\tr(LY)$, where
$\bra xL\ket x\geq0$ for every $\ket x$, hence $L\geq0$.
Idempotence then gives
\begin{align}
	\mc E^2(Y)
	        & =R\tr(LR)\tr(LY)=\mc E(Y),
	        &
	\tr(LR) & =1.
	\label{eq:weight-normalization}
\end{align}
The replacement $(R,L)\mapsto(tR,t^{-1}L)$ leaves the map unchanged.
It is only an ambiguity in its factorization.

\subsection{Rank-deficient weights cannot give span five}

The output--input Choi matrix of \cref{eq:rank-one-map} is
\begin{align}
	J(\mc E)
	 & =\sum_{v,v'}\mc E(\ket v\bra{v'})
	\otimes\ket v\bra{v'}
	\nonumber                                    \\
	 & =\sum_{v,v'}RL_{v'v}\otimes\ket v\bra{v'}
	=R\otimes L^{\mathsf T}.
	\label{eq:transfer-choi}
\end{align}
It is also the sum of the outer products of the four vectorized
Kraus matrices. Every Kraus vector lies in its support: if a vector
is in the kernel of this positive sum, the sum of its squared inner
products with all four Kraus vectors is zero, so each inner product
is zero.

If $\rank R=1$, every Kraus matrix has a common one-dimensional
output range and can be written as $\ket r\bra w$.
Every product $K\otimes\overline{K'}$ then has range in the fixed
line spanned by $\ket r\otimes\overline{\ket r}$. The space of maps
from $\C^4$ into a fixed line has dimension four.
If $\rank L=1$, every Kraus matrix has a common one-dimensional
input support. Their tensor products vanish on a fixed
three-dimensional subspace and again span at most four dimensions.
All physical coordinates in \cref{eq:letters} are linear
combinations of these products.

\begin{proposition}
	\label{prop:faithful}
	For a rank-one trace transfer, full coefficient span five requires
	$R>0$ and $L>0$. Every rank-deficient input or output weight is
	excluded, including nonminimal four-Kraus representations of such
	a transfer.
\end{proposition}

The same argument explains the prior environment-three exclusion
in the accepted specification~\cite{Targets}. With three trace
Kraus matrices, the Choi rank obeys
$\rank R\,\rank L\leq3$. At least one qubit weight must then have
rank one, and coefficient span five is impossible. Environment four
is the first case in which both weights can be faithful.

\subsection{Exhaustive parameterization of the faithful locus}

Apply \cref{eq:virtual-gauge} with $S=L^{1/2}$. The transformed map is
\begin{align}
	\mc E'(Y)
	 & =L^{1/2}\mc E(L^{-1/2}YL^{-1/2})L^{1/2}
	\nonumber                                  \\
	 & =L^{1/2}RL^{1/2}
	\tr(LL^{-1/2}YL^{-1/2})
	\nonumber                                  \\
	 & =L^{1/2}RL^{1/2}\tr Y.
	\label{eq:whitening}
\end{align}
The new output weight has trace $\tr(LR)=1$. A unitary virtual
similarity diagonalizes it. Ordering its eigenvalues gives
\begin{align}
	L=\Id_2,\qquad
	R=\diag(p,1-p),\qquad
	0<p\leq1/2.
	\label{eq:normal-form}
\end{align}
This is a virtual gauge reduction of the given tensor. No scalar
rescaling of the physical tensor has occurred.

Write $p_0=p$ and $p_1=1-p$. The matrix units
$E_{mn}=\ket m\bra n$ give a reference Kraus family
$K_{mn}=\sqrt{p_m}E_{mn}$, since
\begin{align}
	\sum_{m,n}K_{mn}YK_{mn}^\dagger
	 & =\sum_{m,n}p_mY_{nn}\ket m\bra m
	\nonumber                           \\
	 & =R(Y_{00}+Y_{11})=R\tr Y.
	\label{eq:reference-kraus}
\end{align}
Its Choi matrix $R\otimes\Id_2$ has rank four. Every representation
by four Kraus matrices is therefore minimal and consists of four
linearly independent matrices.

The unitary freedom of these representations follows directly from
factorization. Let $Z$ and $Z_0$ have the vectorized Kraus matrices
of an arbitrary family and the reference family as columns.
Both are invertible and
\begin{align}
	ZZ^\dagger & =Z_0Z_0^\dagger,
	           &
	U          & =Z_0^{-1}Z,
	\nonumber                                          \\
	UU^\dagger
	           & =Z_0^{-1}ZZ^\dagger(Z_0^\dagger)^{-1}
	=\Id_4.
\end{align}
Thus $U$ is unitary. Equivalently, all tensors in the faithful
stratum have representatives
\begin{align}
	A_0 & =\sqrt R\,H_0,
	    &
	A_1 & =\sqrt R\,H_1,
	    &
	C_0 & =\sqrt R\,H_2,
	    &
	C_1 & =\sqrt R\,H_3,
	\label{eq:frame-param}
\end{align}
where $H_0,H_1,H_2,H_3$ are a Hilbert--Schmidt orthonormal basis of
$M_2(\C)$. Their defining conditions and completeness relation are
\begin{align}
	\tr(H_a^\dagger H_b)
	 & =\sum_{u,v}\overline{(H_a)_{uv}}(H_b)_{uv}
	=\delta_{ab},
	\nonumber                                     \\
	\sum_{a=0}^3(H_a)_{uv}\overline{(H_a)_{u'v'}}
	 & =\delta_{uu'}\delta_{vv'}.
	\label{eq:frame-equations}
\end{align}

Modulo the active physical unitary and the scalar Kraus unitary,
the frame data reduce to the two-plane
$\mc V=\operatorname{span}\{H_0,H_1\}$.
Its orthogonal complement determines the scalar coordinates.
The reduced parameter space is therefore
\begin{align}
	\left\{(p,\mc V):
	0<p\leq1/2,\quad
	\mc V\in\operatorname{Gr}(2,M_2(\C))\right\}/\mc G_p.
	\label{eq:parameter-space}
\end{align}
Here $\operatorname{Gr}(2,M_2(\C))$ denotes the set of complex
two-dimensional subspaces, and $\mc G_p$ consists of unitary
conjugations preserving $R$. For $p<1/2$ these are diagonal
unitaries in the eigenbasis of $R$; for $p=1/2$ they are all
qubit unitaries. Central phases act trivially.

There is no additional nonunitary residual virtual similarity in
this normalization. Under a similarity $S$, the input weight becomes
$(SS^\dagger)^{-1}$. It can be returned to $\Id_2$ by the scalar
factorization ambiguity only if $SS^\dagger$ is scalar. The scalar
factor in $S$ cancels from $SKS^{-1}$, leaving a unitary similarity.
Thus \cref{eq:parameter-space} is exhaustive modulo precisely the
gauges stated here. It does not assert a classification under
arbitrary physical transformations.

\subsection{Coefficient span five is automatic after reduction}

Define the horizontal coordinate matrices
\begin{align}
	P_{jk} & =A_j\otimes\overline{A_k},
	       &
	Q      & =C_0\otimes\overline{C_0}
	+C_1\otimes\overline{C_1}.
	\label{eq:coordinates}
\end{align}
The four $P_{jk}$ are independent because $A_0,A_1$ are independent.
Indeed, dual linear functionals on their two-dimensional span
extract the coefficient of each tensor product.

The transfer $T$ is not in their span. Define the invertible
entry-rearrangement map $\mc R$ by
\begin{align}
	\mc R(K\otimes\overline{K'})
	 & =\operatorname{vec}(K)\operatorname{vec}(K')^\dagger.
\end{align}
Then
\begin{align}
	\mc R(P_{jk})
	         & =\operatorname{vec}(A_j)\operatorname{vec}(A_k)^\dagger,
	         &
	\mc R(T) & =R\otimes\Id_2.
	\label{eq:reshuffle}
\end{align}
Every linear combination of the first four matrices has range in
$\operatorname{span}\{\operatorname{vec}(A_0),
	\operatorname{vec}(A_1)\}$ and hence rank at most two.
The last matrix has rank four. Finally,
\begin{align}
	Q & =T-P_{00}-P_{11}.
	\label{eq:Q}
\end{align}
Thus $P_{00},P_{01},P_{10},P_{11},Q$ are independent.

Every point of \cref{eq:parameter-space} consequently has Gram ranks
$(1,2)$ and coefficient span five. The rank assumptions have not
been replaced by a genericity condition: they hold throughout the
faithful reduced domain.
